\bookchapter{נספח E: תודות והבהרות}

תודה שהתמדתם וקראתם עד הסוף את “קורס הסלולר” הזה! במהלך הכתיבה הסתמכנו על מדריכים רשמיים של יצרניות רבות, על מדריכי משתמש ועל מאמרי הסבר טכנולוגיים פומביים, כדי להבטיח את דיוק המידע.

הטכנולוגיה מתקדמת בקצב מסחרר, ותוכן הספר מבוסס על הטכנולוגיות והמוצרים המרכזיים בעת הכתיבה. כשתגיעו אליו, ייתכן שכבר יופיעו פיצ׳רים ומוצרים חדשים. אנו מקווים שהגישה והשיטות שהספר מלמד יעזרו לכם להתמודד בראש שקט גם עם גאדג׳טים דיגיטליים חדשים בעתיד.

\booksection{עדכוני גרסה 1.6}

\begin{itemize}
  \item הרחבת הנושא "ערך נכסים דיגיטליים ועקרונות ניהול", והוספת פרק "ניהול תמונות ומערכת קבצים" (מתן שמות לפי זמן+אירוע, מיון וארגון, ניקוי אוטומטי אחרי 72 שעות, הסרת נתוני EXIF).
  \item שכלול מערכת ההגנה של "מצב נקי/הגנה" וההגנה מפני "טפילים דיגיטליים".
  \item הוספת תרתי משמע הונאה נפוצים בעידן ה-AI (החלפת פנים וחיקוי קול / התחזות לגורם סמכותי / השקעות פיקטיביות) וגישות הגנה להתמודדות.
  \item שכתוב ההגנה מפני דוקסינג ("פתיחת קופסה") והתקפות הנדסה חברתית.
  \item הרחבת פרק הנגישות (הפעלה מלאה של VoiceOver/TalkBack, הגדלה ומסנני צבע, תיאור תמונות ב-AI, כתוביות חיות, ושני ערוצי שיחה RTT/TTY).
  \item שכתוב מסגרת הפרומפטים לפי הגישה "תפקיד-הקשר-משימה-הגבלות", עם 3 מקרי בוחן שלמים.
  \item הוספת שלושה נספחים טכנולוגיים למתקדמים: נספח לפרק 1: ניתוח הרחבת זיכרון בטלפוני אנדרואיד (ההגדרה העצמית, עקרונות zRAM/swap, צווארי הבקבוק במהירות, גבולות התרחישים הרלוונטיים, ושמות ואופני הפעלה אצל היצרניות השונות).
  \item נספח לפרק 20: Termux — סביבת מסוף Linux ניידת על מכשירי Android (פונקציות הליבה, אקוסיסטם התוספים, התקנה ותצורה, ותרחישים אופייניים כמו פיתוח נייד, תחזוקת שרתים ובדיקות אבטחה).
  \item נספח לפרק 23: מדריך למתחילים ל-GitHub — מאחסון קוד ועד פלטפורמת שיתוף פעולה (ההבחנה בין Git ל-GitHub, מושגי ליבה, הרשמה והתחלה, זרימת שיתוף פעולה עם Pull Request, תזכורות אבטחה, ופתרונות גישה מסין).
\end{itemize}

\booksection{עדכוני גרסה 1.7}

א. הוספת 4 פרקים חדשים

\begin{itemize}
  \item פרק 29: בריאות דיגיטלית — ניתוח זמן שימוש במסך, מגבלות אפליקציות, מצב נא לא להפריע, כלל 20-20-20 לשמירת העיניים, הצלת עמוד הצוואר, ניהול שינה, "עיוורון סלקטיבי" מול חרדת מידע, ואתגר בריאות דיגיטלית של שבוע.
  \item פרק 30: שיקוף מסך ורכב — השוואת שלושת המחנות AirPlay / Miracast / DLNA, פתרונות שיקוף בכבל, מצב שולחן עבודה Samsung DeX, ארבעת פתרונות הרכב הגדולים CarPlay / Android Auto / HiCar / CarLife, ושיקוף מסך מעשי בחדרי ישיבות.
  \item פרק 31: יצירת סרטונים קצרים — שלוש שאלות לבחירת נושא, הגדרות צילום ב-5 שלבים, חמש טכניקות תנועת מצלמה (התקרבות, התרחקות, סיבוב, סחיבה ומעקב), זרימת עבודה בת שש השלבים ב-CapCut, יסודות הכנת תמונת שער, ותורת התוכן של "3 השניות הראשונות קובעות הכל".
  \item פרק 32: מדריך הצלה ואבטחה — נוהל 30 הדקות המוזהבות לאובדן טלפון, תהליך 72 השעות לטיפול אחרי הונאה, איתור בעיות במצב בטוח אחרי זיהום, שישה צעדים להשבת חשבון שנגנב, התמודדות עם הודעות סחיטה, וחמש פעולות מנע יומיומיות.
\end{itemize}

ב. הוספת שני פרקים חדשים

\begin{itemize}
  \item פרק 27: התאמה לגיל השלישי וקשרי משפחה — טבלת השוואה של מצבים פשוטים לפי יצרניות, שש מלכודות הונאה מרכזיות ושלוש חזיתות הגנה, סיוע מרחוק (MeeTime / Sunflower), מצב האכפתיות של WeChat וקלט קולי, ו-SOS / זיהוי נפילה / תזכורות תרופות.
  \item פרק 28: מדריך רכישה — שיטת בחירת טלפון בשלושה צעדים, טבלת מפרט "רק חמישה פריטים", המלצות לפי חמש רמות תקציב, טבלת בחירה לשבעה תרחישים, חמש המלכודות הגדולות ברכישה, וניתוח תזמון הקנייה וערוצי הרכישה.
\end{itemize}

ג. הוספת תכונות הקראה ואינטראקציה

\begin{itemize}
  \item סרגל בקרת הקראה צף: תמיכה בנגינה/השהיה, מקטע קודם/בא, התאמת מהירות 0.5x–2.5x, ועצירה.
  \item כפתור אוזניות מהיר בצד ימין: הפעלת הקראה בלחיצה אחת, עם התחלה אוטומטית ממקום הקריאה הנוכחי.
  \item קיצורי מקלדת: R להפעלה/כיבוי של ההקראה, ESC לעצירה.
  \item שיפור חוויית ההקראה: מעבר רציף ואוטומטי בין פרקים, והסתרה אוטומטית של פקדי ההקראה בהדפסה.
\end{itemize}

ד. עדכוני מערכת

ריענון מלא ומסונכרן של תוכן העניינים, מסלול הלמידה (מארבעה שלבים לחמישה), מפת הפרקים, מספר הגרסה (מ-1.6 ל-1.7), מספור הנספחים (מ-5 ל-6) והמטא-דאטה.

ה. עדכוני עדכניות התוכן

\begin{itemize}
  \item תקני מעבדים: Snapdragon 8 Elite Gen5 / Dimensity 9500 מסווגים ברמה העליונה (בסיס 2026).
  \item דגמי דגל להתייחסות: iPhone 17 Pro, Samsung S26 Ultra, Xiaomi 17 Ultra, Huawei Pura 80 Ultra, vivo X300 Pro.
  \item פרקי AI: עודכנו ל-DeepSeek R1, ChatGPT-5, ותקן כוח החישוב 80 TOPS של NPU במכשיר עצמו.
  \item מסכים מתקפלים: טווח המחירים ירד לרמה של 4,000–5,000 יואן.
  \item מדריך רכישה: כל רמות התקציב וההמלצות לפי תרחישים עודכנו לשוק של 2026.
  \item רשת ואקוסיסטם: Wi-Fi 7 / 5G-A / אקוסיסטם HarmonyOS / LE Audio — הכל עודכן לתקנים העדכניים.
\end{itemize}

ו. מודול תכונות חדש: לוח מחווני למידה (מתחת לאזור ה-Hero)

\begin{itemize}
  \item גרף טבעתי להתקדמות קריאה: מציג בזמן אמת את אחוז הפרקים שנקראו, עם אנימציית טבעת SVG.
  \item בדיקה עצמית לבריאות הטלפון: רשימה של 8 בדיקות אבטחה (סיסמה, אימות דו-שלבי, איתור הטלפון, גיבוי ענן, גיבוי WeChat, אבטחת תשלומים, אפליקציית אנטי-הונאה, ניקוי אישורי תשלום ללא סיסמה).
  \item שמירת מצב: מצב הסימונים נשמר ב-localStorage, ציון הבריאות מחושב דינמית ומוצג בפס התקדמות מדורג.
  \item הודעות בדירוגים: משוב שונה לפי הציון — "ציון מלא / טוב / עובר / סיכון אבטחה".
\end{itemize}

\booksection{עדכוני גרסה 1.8}

הוספת החלק השביעי “פרקטיקת תכנות בעזרת AI” (פרקים 37–47, 11 פרקים בסך הכל): המהפכה הפרדיגמטית ו-Agentic Engineering, חמשת עקרונות החשיבה של תכנות בעזרת AI, הנדסת הקשר ופיתוח מונחה מפרט (SDD), הנדסת דרישות, פנורמת כלי התכנות בעזרת AI ל-2026, פנורמת המודלים הגדולים ועבודה מעשית עם API, הנדסת שיחה, פרויקט שלם מאפס לאחת בזרימה מלאה, הבטחת איכות, דרך ההתקדמות (MCP/שיתוף פעולה בין סוכנים מרובים) ומדריך עיון מהיר. ריענון מסונכרן של תוכן העניינים, מסלול הלמידה (מחמישה שלבים לשישה), מפת הפרקים ומדד העיון המהיר. התוכן מבוסס על מהדורת אוגוסט 2026 של “חשיבת תכנות בעזרת AI: מדריך מתודולוגי מאפס ועד פרקטיקה מלאה”.

הוספת החלק השמיני “פרקטיקת פיתוח אנדרואיד” (פרקים 48–58, 11 פרקים) והחלק התשיעי “פרקטיקת פיתוח iOS” (פרקים 59–69, 11 פרקים): שתי הפלטפורמות מכסות הקמת סביבת פיתוח (Android Studio / Xcode), לימוד מזורז של השפה (Kotlin / Swift), פריימוורק ממשק (Jetpack Compose / SwiftUI), נתונים ורשת, שני פרויקטים מעשיים שלמים (ספר חשבונות + חיפוש מזג אוויר), בדיקות ועלייה לחנות (Google Play / App Store) ומדריך עיון מהיר. ריענון מסונכרן של תוכן העניינים, מסלול הלמידה (משישה שלבים לשמונה), מפת הפרקים ומדד העיון המהיר. התוכן מבוסס על מהדורות אוגוסט 2026 של “פיתוח אנדרואיד מאפס” ו“פיתוח iOS מאפס”.

\booksection{השלמות תוכן: תוכנה שימושית וחסימת ספאם והטרדות}

\begin{itemize}
  \item בפרק 13 נוסף סעיף שישי "כלי מערכת שימושיים": כלי ניקוי כמו Qingzhuo, ניקוי האחסון המובנה של המערכת, המרכז הלאומי למניעת הונאות, אפליקציות ניהול הטלפון של היצרניות, מניעת הטרדות של הספקים ו-AdGuard — יחד עם עקרונות הבחירה "מערכת קודם / השלמה רשמית / ריסון בצד שלישי".
  \item טבלת השוואה של חסימת ספאם מובנית לפי יצרניות: נתיבי כניסה והגדרות מומלצות עבור iOS, Xiaomi/Redmi, Huawei/Honor, OPPO/realme/OnePlus, vivo/iQOO ו-Samsung.
  \item הרחבת יכולות חסימת הספאם בפרק 8: הוספת הסברים מעשיים על זיהוי מתקשרים, רשימות שחורות ולבנות, סינון SMS לפי מילות מפתח, השתקת מתקשרים לא מוכרים ב-iOS ותיאום עם מצב נא לא להפריע — עם קישורים הדדיים לפרק 13.
  \item המלצה על הגנה תלת-שכבתית: חסימה במערכת + מניעת הטרדות של הספק + המרכז הלאומי למניעת הונאות, תוך הדגשת חשיבות רשימה לבנה למניעת חסימות שגויות.
\end{itemize}

\booksection{השלמות תוכן: התרעות מזג אוויר / התרעות אסון טבע / התרעות רעידות אדמה}

\begin{itemize}
  \item בפרק 13 נוסף סעיף שביעי: התרעות מזג אוויר (ערוצי מטאורולוגיה רשמיים + אפליקציית מזג האוויר של המערכת + דירוגי צבע), התרעות אסון טבע (ערוצים מרובים: מטאורולוגיים/חירום), ועקרונות התרעות רעידות אדמה וטעויות נפוצות (התרעה $\neq$ תחזית).
  \item טבלת הפעלת התרעות רעידת אדמה לפי יצרנית: הכניסות הנפוצות ונקודות לשים לב בהרשאות עבור Xiaomi, Huawei/Honor, סדרת OPPO, vivo, iPhone ו-Samsung.
  \item בפרק 32 נוסף "תרחיש אפס": סדר פעולות ההגנה אחרי אזעקת רעידת אדמה, התמודדות עם גשמים עזים/טייפון, והחיבור ל-SOS לווייני / תקשורת ללא רשת; מדריך ההצלה מקושר להגדרות ההתרעות בפרק 13.
  \item תיאום עם סעיף מזג האוויר: הערך על אפליקציות מזג אוויר בפרק הנסיעות הושלם בהתראות אזהרה ובתזכורת "ההתראות מגיעות רק אם הפעלתם אותן".
\end{itemize}

\booksection{השלמות תוכן: שיתוף פעולה בין מכשירים וכלי פרודוקטיביות}

\begin{itemize}
  \item סעיף שמיני בפרק 13: שליטה במחשב מהטלפון (Sunflower, ToDesk) והקווים האדומים של אנטי-הונאה; משרד עם AI (Trae / WorkBuddy / Dumate) וזרימת עבודה לשליטה במחשב מהטלפון; טפט דינמי מבית miHoYo (HoYo Play).
  \item העברה מהירה ברשת מקומית ולוח גזירה: LocalSend, Feige Chuanshu, AirDrop; KDE Connect, לוח גזירה משותף, והעברה בין טלפון למחשב של היצרניות השונות.
  \item קודים דינמיים לאימות דו-שלבי (2FA): Aegis, Microsoft Authenticator, קודי גיבוי ונקודות לתשומת לב בהחלפת מכשיר.
  \item TickTick וסריקת מסמכים נייר: הרחבת השימושים ללוח שנה ולמשימות; השלמת כלי סריקה וטכניקות צילום.
\end{itemize}

מאחלים לך לשלוט בטלפון בקלות, ליהנות מהנוחות וההנאה שהטכנולוגיה מביאה — ולהפוך לאזרח דיגיטלי ערני ובטוח בעצמו.

— המדריך שלך לחיים הדיגיטליים
